% Vietnamese translation for OI Public Library
% Source-PDF: 英文资料 ENGLISH MATERIALS/Introduction To Algorithm - 3rd Edition.pdf
\documentclass[11pt]{article}
\usepackage{oipl-vn}

\title{Nhập môn thuật toán}
\author{Thomas H. Cormen, Charles E. Leiserson, Ronald L. Rivest, Clifford Stein}
\date{MIT Press, ấn bản thứ ba, 2009}

\begin{document}
\maketitle

\origtitle{Introduction to Algorithms, Third Edition}
\sourcepath{English Materials / Introduction To Algorithm - 3rd Edition.pdf}

\renewcommand{\abstractname}{Tóm tắt}
\begin{abstract}
Đây là ghi chú dịch tiếng Việt cho \emph{Introduction to Algorithms}, ấn bản
thứ ba. Sách là tài liệu tham khảo kinh điển về thuật toán và cấu trúc dữ liệu,
bao phủ nền tảng, sắp xếp, cấu trúc dữ liệu, thiết kế thuật toán, đồ thị, số
học, xâu, hình học tính toán, NP-đầy đủ và thuật toán xấp xỉ. Ghi chú này
chuyển ngữ thông tin đầu sách, cấu trúc phần và các chủ đề chính.
\end{abstract}

\section{Thông tin sách}

\begin{center}
\begin{tabular}{ll}
\textbf{Tác giả} & Thomas H. Cormen; Charles E. Leiserson; Ronald L. Rivest; Clifford Stein \\
\textbf{Nhà xuất bản} & The MIT Press \\
\textbf{Ấn bản} & Thứ ba \\
\textbf{Năm} & 2009 \\
\textbf{ISBN} & 978-0-262-03384-8
\end{tabular}
\end{center}

Sách được dùng rộng rãi trong các khóa thuật toán bậc đại học và sau đại học.
Mỗi chương tương đối độc lập, trình bày thuật toán, kỹ thuật thiết kế, phân
tích độ phức tạp và bài tập. Các đoạn mã giả trong sách là tài liệu tham khảo,
không phải code template cần dịch lại.

\section{Cấu trúc phần đã dịch}

\begin{description}
  \item[Phần I. Nền tảng]
  Vai trò của thuật toán, sắp xếp chèn, phân tích độ phức tạp, ký hiệu tiệm
  cận, chia để trị, giải truy hồi và thuật toán ngẫu nhiên.

  \item[Phần II. Sắp xếp và thống kê thứ tự]
  Heapsort, quicksort, sắp xếp tuyến tính, chọn phần tử thứ \(k\), median và
  các cận dưới cho bài toán sắp xếp.

  \item[Phần III. Cấu trúc dữ liệu]
  Stack, queue, linked list, hash table, cây tìm kiếm nhị phân, cây đỏ--đen và
  kỹ thuật tăng cường cấu trúc dữ liệu.

  \item[Phần IV. Kỹ thuật thiết kế và phân tích nâng cao]
  Quy hoạch động, thuật toán tham lam, matroid, mã Huffman và phân tích khấu
  hao.

  \item[Phần V. Cấu trúc dữ liệu nâng cao]
  B-tree, Fibonacci heap, van Emde Boas tree và cấu trúc tập rời nhau với nén
  đường đi.

  \item[Phần VI. Thuật toán đồ thị]
  BFS, DFS, sắp xếp topo, thành phần liên thông mạnh, cây khung nhỏ nhất,
  đường đi ngắn nhất một nguồn, mọi cặp đường đi ngắn nhất và luồng cực đại.

  \item[Phần VII. Chủ đề chọn lọc]
  Thuật toán đa luồng, phép toán ma trận, quy hoạch tuyến tính, FFT, thuật toán
  số học, khớp xâu, hình học tính toán, NP-đầy đủ và thuật toán xấp xỉ.

  \item[Phần VIII. Phụ lục nền tảng toán học]
  Tổng, tập hợp, quan hệ, hàm, đồ thị, cây, đếm, xác suất rời rạc và ma trận.
\end{description}

\section{Ghi chú nội dung}

Trong luyện OI, các chương thường dùng nhất là:
\begin{itemize}
  \item Chương 1--5 cho nền tảng phân tích thuật toán, chia để trị và ngẫu nhiên.
  \item Chương 6--9 cho sắp xếp, heap, quicksort, radix/counting sort và chọn
        thống kê thứ tự.
  \item Chương 10--21 cho cấu trúc dữ liệu cơ bản và nâng cao.
  \item Chương 22--26 cho thuật toán đồ thị kinh điển.
  \item Chương 30--33 cho FFT, số học, khớp xâu và hình học tính toán.
  \item Chương 34--35 cho NP-đầy đủ và xấp xỉ, hữu ích để nhận diện giới hạn
        của các bài toán.
\end{itemize}

\section{Ghi chú về trích xuất}

Tệp PDF có 1313 trang và lớp văn bản trích xuất tốt. Một số ký hiệu trong mục
lục, như dấu đánh sao cho phần nâng cao, có thể hiện thành ký tự lạ khi dùng
\texttt{pdftotext}. Bản dịch đầy đủ từng chương cần đối chiếu PDF gốc khi
chuyển công thức, hình, bảng và mã giả sang LaTeX.

\end{document}
